\subsection{The common raw receiver}
\label{sec:conclusion}

We now use the coefficient cover and regional raw budgets to define one
interval on which all regional estimates hold for the complete return.

\begin{lemma}[The common raw receiver]\label{lem:common-raw-receiver}
Fix the coefficient cover and reserves of
Lemma~\ref{lem:pure-coefficient-cover}, and let \(b>0\) be the raw
budget of Proposition~\ref{prop:regional-raw-budgets}. For every fixed
\(0<h_c\le b\), there are a raw coefficient neighborhood and a
fixed \(\bar m>0\) such that all covering regional constructions
and their every-point estimates hold on the same actual interval
\begin{equation}\label{eq:common-raw-carrier}
 I_\theta(h_c)=\{x:a_\theta+x\in J_\theta^{\rm out},\quad
                    0<h_\theta(a_\theta+x)<h_c\}.
\end{equation}
This interval retains the original input and inverse output windows,
positive-time passage, speed and stop clips, admitted equality dividers,
and all intervening nonroot inputs. The normalized regional descriptions
use their exact centres and positive absolute gains and the same ordinary
\(C^3\) divided closing. The coefficients and actual \(m\) are fixed
in every phase derivative. The neighborhood and multiplier bound may
depend on \(h_c\), but the coefficient cover and its reserves do not.
\end{lemma}

\begin{proof}
Fix any \(0<h_c\le b\). Proposition~\ref{prop:regional-raw-budgets}
gives the finitely many independent bounds
\(\delta_j(h_c),\bar m_j(h_c)\) on the full actual regional
domains. Intersect their raw neighborhoods, including
\(|\alpha_0|<d_{\rm out}/4\). Let \(m_{\rm ad}>0\) be the
ordinary closing-domain bound, and choose
\begin{equation}\label{eq:common-multiplier-bound}
\bar m=\tfrac12\min\{m_{\rm ad},\bar m_1(h_c),\ldots,
                          \bar m_N(h_c),(2C_H)^{-1}\}>0.
\end{equation}
The complete displacement \eqref{eq:return} is therefore defined
on every admitted regional image, with its original divided observer
\eqref{deg:det-divided-closing}. All three nonlinear curvature terms
in \eqref{eq:raw-budget-complete-curvature} remain in its estimates.

The original raw admission set is exactly
\[
J_\theta^{\rm out}=
\mathcal T_\theta\cap\operatorname{int}K_+^\sharp
\cap T_\theta^{-1}(\operatorname{int}K_-^\sharp)
\cap\{\sigma_\theta<d_{\rm out}\}.
\]
Scalar uniqueness and first variation order the actual through graphs.
Thus \(\mathcal T_\theta\), both window preimages, the speed clip,
and the central-height clip are intervals; central height is strictly
increasing by \eqref{loc:central-height}. Their intersection
\eqref{eq:common-raw-carrier} is one interval, possibly empty.

The regional lower, core and high sets cover each admitted through
input before clipping. Intersect them with this same interval. Their
actual matching identities, inverse images and every-point estimates
survive. In particular \eqref{deg:det-height-cover},
\eqref{deg:trace-domain-inclusion}, \eqref{deg:negative-gamma-cover},
\eqref{deg:far-physical-cover}, \eqref{deg:compact-physical-cover},
and Proposition~\ref{p2:prop-unified-carrier-cover} retain their
admitted positive equality dividers. The compact positive-ratio sets
remain the monotone height preimages of
Lemma~\ref{ann:lem-compact-positive-domains}. If a clip removes an
overlap, the retained connected interval is still covered by the
adjoining enlarged estimates. The natural stopped endpoints are
excluded. The original compact-cubic and bounded-positive ghost raw
tails retain \eqref{deg:compact-raw-tail-derivative} and
\eqref{bp-tail:consumed-derivative}, respectively.

The fixed cap is independent of the normalized overlap heights.
Proposition~\ref{prop:regional-raw-budgets} places all comparison
shots and upper core endpoints below \(h_c/4\) by later raw
restrictions and covers the high complement up to \(h_c\).
For degree loss it uses
\eqref{eq:degree-cap-height-raw-bound}--\eqref{eq:degree-cap-alpha-one-bound}
on both the compact and reciprocal rectangles. Hence no ratio reserve
or coefficient threshold is selected again when \(h_c\) is reduced.
Set \(h_{\rm out}=h_c\) in \eqref{loc:analytic-receiver}.

Finally \eqref{eq:finite-multiplier}, with
\(\eps_f=h(\theta)K\) and positive \(h\), gives the actual
\(m(\theta)\le\bar m\) on \(K>0\) after a further raw coefficient
shrink. This restriction is independent of \(s,|\gamma|\) and of
all root gaps. For \(K\le0\) the selected finite channel has no
positive-time exit. Thus the receiver and all its regional estimates
are fixed before any zero count or geometric capture.
\end{proof}

One admissible fixed choice for the subsequent argument is
\begin{equation}\label{eq:common-raw-cap}
h_c=b/2>0.
\end{equation}
